\documentclass[11pt,letterpaper]{article}
\usepackage[margin=1in]{geometry}
\usepackage{booktabs}
\usepackage{amsmath}
\usepackage{amssymb}
\usepackage{hyperref}
\usepackage{float}
\usepackage{caption}
\usepackage{parskip}
\usepackage{listings}
\usepackage{xcolor}
\setlength{\parskip}{6pt}

\lstset{
  basicstyle=\ttfamily\footnotesize,
  breaklines=true,
  frame=single,
  keepspaces=true,
  showstringspaces=false,
  language=C++,
  commentstyle=\color{gray}\itshape,
  keywordstyle=\color{blue!60!black}\bfseries,
}

\title{\textbf{Adaptive Stream Buffer Prefetching:\\
Implementation, Bug Hunting, and Empirical Evaluation}}
\author{Aengus McGuinness \and Ali Saffrini \and Madison Gates \and Jacob Tom\\[0.5em]
\textit{CS 1411 --- Spring 2026 Final Project}}
\date{}

\begin{document}
\maketitle

%─────────────────────────────────────────────────────────────────────────────
\begin{abstract}
%─────────────────────────────────────────────────────────────────────────────
We present an empirical evaluation of two stream-buffer prefetchers built on a
shared Pin-based simulator: (1)~a fixed-depth FIFO stream buffer following
Jouppi~\cite{jouppi1990improving}, and (2)~an adaptive stream buffer with a
two-level cache hierarchy and a learned-histogram depth policy following
Palacharla and Kessler~\cite{palacharla1994evaluating}. We sweep both
implementations across three SPEC~CPU2006 workloads
(\texttt{libquantum}, \texttt{hmmer}, \texttt{dealII}). Phase~1 establishes
strong workload sensitivity: prefetching reduces effective demand misses by
$\sim$43$\times$ on \texttt{libquantum} but bloats L2 traffic by up to
9$\times$ on irregular workloads with negligible payoff. Phase~2 reproduces
this picture on a richer cycle model and introduces an adaptive depth knob:
at maximum depth, adaptive achieves up to $5.02\times$ speedup on
\texttt{libquantum} (vs.\ next-line's $5.83\times$); at depth~$2$, adaptive
attains higher prefetch accuracy than next-line on \texttt{dealII}
($83.3\%$ vs.\ $78.4\%$) while preserving comparable speedup. We also document
three implementation bugs in the adaptive policy whose correction was
necessary to recover the paper's claimed performance, raising adaptive's
\texttt{libquantum} speedup from $1.01\times$ to $5.02\times$.
\end{abstract}

%─────────────────────────────────────────────────────────────────────────────
\section{Introduction}
%─────────────────────────────────────────────────────────────────────────────
Memory latency dominates many computational kernels and grows worse as core
counts and clock rates scale faster than DRAM. Hardware prefetching is a
classical mitigation: speculatively issue memory requests before the demand
access occurs. The simplest variant is the \emph{stream buffer}, introduced by
Jouppi~\cite{jouppi1990improving}: a small FIFO of consecutive lines maintained
just outside L1 that intercepts misses before they reach L2.

Naive stream buffers prefetch a fixed number of lines per stream regardless of
program behavior, which wastes memory bandwidth on workloads whose access
patterns are not sequential. Palacharla and
Kessler~\cite{palacharla1994evaluating} address this with an \emph{adaptive}
policy: the prefetcher tracks observed stream lengths in a histogram and uses
that distribution to decide how many lines to prefetch. The histogram lets the
prefetcher behave aggressively on streaming workloads and conservatively on
irregular ones.

This project implements both designs in a shared simulator and compares them
on three SPEC~CPU2006 benchmarks. Section~\ref{sec:phase1} describes the
fixed-depth stream buffer (Phase~1). Section~\ref{sec:phase2} describes the
adaptive stream buffer with an L1D/L2 cache hierarchy (Phase~2), including
three implementation bugs that we identified and fixed.
Section~\ref{sec:methodology} details the sweep methodology, and
Section~\ref{sec:results} reports the headline data and discusses the
comparison.

%─────────────────────────────────────────────────────────────────────────────
\section{Phase 1: Fixed-Depth Stream Buffer}
\label{sec:phase1}
%─────────────────────────────────────────────────────────────────────────────

\subsection{Design}
Phase 1 (\texttt{aengus/stream.cpp}) is a single-translation-unit Pintool that
models the hierarchy
$$\text{L1D} \rightarrow \text{stream buffer} \rightarrow \text{L2} \rightarrow \text{memory}.$$
Each of $S$ parallel streams holds up to $D$ prefetched lines in a strict
FIFO. On every L1D miss, the simulator probes the head of every active
stream:
\begin{itemize}
  \item \textbf{Head match:} the line is popped, served as the demand
        access, and a fresh prefetch is issued one line past the new tail
        (keeping the FIFO full).
  \item \textbf{Head miss:} a new sequential stream is allocated starting at
        $(\textit{missed\_line}+1)$ and pre-populated with $D$ prefetches.
\end{itemize}
The design has no learning component: depth $D$ is a static knob chosen at
configuration time.

\subsection{Sweep parameters}
The Phase~1 sweep (\texttt{aengus/sweep.sh}) explores depth
$D \in \{0,1,2,4,8\}$ and stream count $S \in \{1,2,4,8\}$, plus an
associativity sweep at $S=4$, $D=4$. Three SPEC~CPU2006 benchmarks are run
natively to completion under Pin instrumentation:
\texttt{libquantum} (quantum simulation, highly streaming),
\texttt{hmmer} (HMM-based sequence search, irregular),
\texttt{dealII} (PDE finite-element solver, mixed).

\subsection{Headline result}
Table~\ref{tab:phase1-eff} reports the number of \emph{effective} L1D demand
misses --- those that survived the stream buffer --- at $S=4$ across all
sweep depths. Table~\ref{tab:phase1-l2} reports the corresponding total L2
request count, which captures the bandwidth cost of the prefetcher.

\begin{table}[H]
\centering
\caption{Phase 1 effective L1D demand misses ($S=4$, lower is better).}
\label{tab:phase1-eff}
\begin{tabular}{rrrr}
\toprule
depth & libquantum & hmmer & dealII \\
\midrule
0 & 121{,}200{,}271 & 1{,}089{,}349{,}815 & 488{,}587{,}529 \\
1 &   2{,}813{,}918 & 1{,}061{,}848{,}629 & 476{,}963{,}724 \\
2 &   2{,}781{,}862 & 1{,}061{,}093{,}631 & 481{,}485{,}540 \\
4 &   2{,}813{,}762 & 1{,}062{,}760{,}925 & 477{,}022{,}769 \\
8 &   2{,}813{,}636 & 1{,}062{,}569{,}741 & 484{,}302{,}817 \\
\bottomrule
\end{tabular}
\end{table}

\begin{table}[H]
\centering
\caption{Phase 1 total L2 requests ($S=4$, demand+prefetch traffic).}
\label{tab:phase1-l2}
\begin{tabular}{rrrr}
\toprule
depth & libquantum & hmmer & dealII \\
\midrule
0 & 121{,}204{,}522 &   1{,}091{,}088{,}059 &   525{,}692{,}727 \\
1 & 124{,}038{,}602 &   2{,}130{,}266{,}731 &   999{,}832{,}393 \\
2 & 126{,}766{,}816 &   3{,}189{,}907{,}621 & 1{,}489{,}849{,}986 \\
4 & 132{,}479{,}525 &   5{,}320{,}405{,}009 & 2{,}430{,}902{,}592 \\
8 & 143{,}733{,}460 &   9{,}570{,}217{,}539 & 4{,}404{,}294{,}345 \\
\bottomrule
\end{tabular}
\end{table}

\subsection{Discussion}
Two patterns emerge:
\begin{enumerate}
  \item \textbf{libquantum collapses by $\sim$43$\times$} at any depth
        $\geq 1$; its access pattern is essentially infinite sequential
        streams, so the very first prefetched line catches almost every
        future miss. Increasing depth beyond 1 yields no further reduction.
  \item \textbf{hmmer and dealII} barely shrink (under 3\%) but their L2
        traffic grows roughly linearly with depth. At $D=8$, \texttt{hmmer}
        issues 9$\times$ the L2 requests of the no-prefetch baseline for a
        $\approx2\%$ miss reduction. \texttt{dealII} pays a similar 8$\times$
        bandwidth tax for a $\approx2\%$ payoff.
\end{enumerate}
A static prefetcher cannot tell these workloads apart; choosing depth
forces a trade-off between \texttt{libquantum}'s win and the
hmmer/dealII bandwidth waste. This sets up the adaptive Phase~2 design.

%─────────────────────────────────────────────────────────────────────────────
\section{Phase 2: Adaptive Stream Buffer}
\label{sec:phase2}
%─────────────────────────────────────────────────────────────────────────────

\subsection{Design overview}
Phase~2 (\texttt{ali/}) extends the Phase~1 model with three additions:
\begin{enumerate}
  \item A two-level cache hierarchy (L1D~$\rightarrow$~L2~$\rightarrow$~memory)
        of independent set-associative LRU caches, ported from the homework~4
        cache model. The stream buffer observes \emph{L2} misses, matching
        Palacharla and Kessler's intent.
  \item A learned-histogram depth policy (next subsection).
  \item A configurable Pin harness with a parallel parameter sweeper that
        streams per-config rows to disk and supports resume across kills.
\end{enumerate}
The codebase splits into a self-contained simulator core
(\texttt{ali/src/stream\_buffer.\{cpp,hpp\}}) reusable by both a standalone
trace-driven CLI and the Pintool, and a Pin instrumentation layer
(\texttt{ali/pintool/adaptive\_stream\_buffer\_pintool.cpp}) that pipes load
and store events into the simulator.

\subsection{Adaptive depth policy}
The simulator maintains up to \texttt{stream\_slots} active streams. Each
slot records the most recently observed line, the inferred direction
($+1$, $-1$, or unknown), and the running stream length. On every L2 miss,
we extend the matching slot or allocate a fresh one.

A cumulative histogram $\textit{lht}_\text{curr}[k]$ counts streams that
\emph{reached length $\geq k$}. After every \texttt{epoch\_reads} L2 misses,
the next epoch's histogram replaces the current one, so the prefetcher
adapts to phase changes. To choose depth on a given access we walk the
histogram from the slot's current length and increment depth while
$$\textit{lht}_\text{curr}[k]<2\cdot\textit{lht}_\text{curr}[k+1],$$
i.e.\ at least half the streams that reached length $k$ also reached
length $k+1$.

\subsection{Three implementation bugs}
\label{sec:bugs}
A naive implementation of the policy above produces results far below
next-line on streaming workloads. We identified and fixed three issues
that were jointly responsible for the under-performance:

\paragraph{Bug A: cliff at the histogram cap.}
The original \texttt{choose\_prefetch\_depth} returned \texttt{0} once a
stream's length reached the histogram bound \texttt{max\_stream\_length},
under the comment ``nothing beyond the largest histogram bucket can be
learned.'' But that condition is exactly when the policy has the
\emph{strongest} evidence the stream is long-lived; truncating to zero
disables prefetching during the most useful phase of the stream's life.
For \texttt{libquantum}'s effectively infinite streams this killed almost
all useful prefetches. Fix: return \texttt{max\_prefetch\_depth} when
\texttt{capped\_length} $\geq$ \texttt{max\_stream\_length}.

\paragraph{Bug B: broken bootstrap.}
Before the first epoch rollover, the policy used a tiny seeded histogram
(\texttt{lht[1]=1}, \texttt{lht[2]=1}, all higher buckets zero). For any
stream of length $\geq 2$ the depth-walking loop hit a zero in the next
bucket on the first probe and returned \texttt{0}, so warmup was almost
entirely silent. Fix: initialize \texttt{history\_ready\_} to \texttt{false}
unconditionally, so the bootstrap path returns \texttt{depth=1} for
\emph{any} length (i.e.\ matches next-line behavior) until a real epoch has
rolled.

\paragraph{Bug D: singleton-stream histogram poisoning.}
When all stream slots were full and a fresh L2 miss could not be tracked,
the original code recorded the orphan miss as a length-$1$ stream
(\texttt{lht\_next\_[1] += 1}). This biased the next epoch's histogram
toward ``streams are short, throttle.'' Fix: \texttt{record\_singleton\_stream()}
is now a no-op; orphan misses contribute no signal.

\paragraph{Quantitative impact.}
Before the fixes, adaptive's best speedup on \texttt{libquantum} was
$1.013\times$ (vs.\ next-line's $5.83\times$) --- effectively no
improvement. After the fixes, adaptive's best speedup rose to
$5.024\times$ (Table~\ref{tab:speedup}), recovering nearly all of
next-line's gain.

\subsection{Cache hierarchy parameterization}
The \texttt{SetAssocCache} class is sized by
$(\textit{size\_bytes}, \textit{line\_size\_bytes}, \textit{associativity})$
with a per-set LRU stack. Defaults follow the homework~4 \texttt{config-base}:
$4$~KB / $1$-way L1D and $1$~MB / $1$-way L2, $64$-byte lines. A read access
probes L1D first, then L2 on miss; an L2 miss is the trigger for the
prefetcher, and the line is installed into both levels on the way back to
the requester.

\subsection{Latency model}
Each access advances logical time by $1$ and adds a per-access cost
\texttt{base\_ref\_cost~+~latency} where \texttt{latency} depends on the
access outcome:

\begin{table}[H]
\centering
\caption{Modeled latency (logical reference steps).}
\begin{tabular}{lr}
\toprule
Outcome & Cost \\
\midrule
L1D hit                                  & 1 \\
L1D miss, L2 hit                         & 10 \\
L2 miss, prefetch ready                  & 1 \\
L2 miss, prefetch in flight              & $\min(80, t_{\text{ready}}-t_{\text{now}})$ \\
L2 miss, no prefetch (full memory miss)  & 80 \\
Write hit (any level)                    & 1 \\
\bottomrule
\end{tabular}
\end{table}
The cycle counts reported in Section~\ref{sec:results} aggregate this per-access
cost over the entire execution.

%─────────────────────────────────────────────────────────────────────────────
\section{Methodology}
\label{sec:methodology}
%─────────────────────────────────────────────────────────────────────────────

\subsection{Workloads}
\begin{itemize}
  \item \textbf{libquantum} (\texttt{libquantum\_O3 400 25}): quantum
        Shor's algorithm simulation. Long arithmetic sweeps over arrays;
        highly sequential.
  \item \textbf{hmmer} (\texttt{hmmer\_O3 nph3.hmm}): HMM-based protein
        sequence search. Pointer-chasing dominated, irregular.
  \item \textbf{dealII} (\texttt{dealII\_O3}): adaptive finite-element
        solver. Mixed access patterns over sparse matrices.
\end{itemize}
All three binaries are compiled at \texttt{-O3} and run inside Pin's
JIT instrumentation. Each Pin invocation is capped at $5\times10^9$
instructions per thread, which we verified is well past the steady-state
phase for all three workloads.

\subsection{Phase 2 sweep parameter space}
The Phase~2 sweep (\texttt{ali/run.sh}) explores
$3 \times 3 \times 2 \times 2 \times 2 = 72$ configurations:

\begin{table}[H]
\centering
\caption{Phase 2 sweep dimensions.}
\begin{tabular}{ll}
\toprule
Dimension & Values \\
\midrule
benchmark             & libquantum, hmmer, dealII \\
policy                & off, nextline, adaptive \\
stream\_slots         & 4, 8 \\
max\_prefetch\_depth  & 2, 8 \\
max\_stream\_length   & 8, 32 \\
\bottomrule
\end{tabular}
\end{table}

The remaining configuration knobs are held at the defaults documented in
\texttt{ali/src/stream\_buffer.hpp} (line size $64$\,B, prefetch buffer $16$
lines, epoch $2000$ misses, stream lifetime $256$ refs, prefetch latency
$8$ refs, miss latency $80$ refs, bootstrap on).

\subsection{Sweep harness}
\texttt{ali/scripts/sweep\_stream\_buffer.sh} runs Pin invocations
$4$-wide in parallel. Each completed configuration writes a single CSV row
to a temporary file, which is appended to the master output CSV
(\texttt{ali/results/stream\_buffer\_experiments.csv}) immediately by the
reaper loop. Two consequences:
\begin{itemize}
  \item Killing the sweep mid-run (cluster time-out, signal, etc.) loses
        at most the four in-flight configurations; everything else is
        already on disk.
  \item Re-running the sweep against the same output CSV skips
        configurations that have an \texttt{ok} row, so resume is a
        single command.
\end{itemize}
Each row records all 22 input dimensions plus the four output metrics
(\textit{selected\_cycles}, \textit{baseline\_cycles}, \textit{speedup},
\textit{accuracy}, \textit{coverage}).

\subsection{Metrics}
\begin{itemize}
  \item \textbf{speedup} $=$ \textit{baseline\_cycles}~$/$~\textit{selected\_cycles},
        where the baseline is an identical simulator run with the prefetch
        policy forced off.
  \item \textbf{accuracy} $=$ \textit{useful\_prefetches}~$/$~\textit{prefetches\_issued}.
        Higher is better; measures how often a prefetch is consumed by a
        demand access.
  \item \textbf{coverage} $=$ \textit{useful\_prefetches}~$/$~\textit{read\_misses}.
        Higher is better; measures the fraction of demand misses that the
        prefetcher rescued.
\end{itemize}

%─────────────────────────────────────────────────────────────────────────────
\section{Results}
\label{sec:results}
%─────────────────────────────────────────────────────────────────────────────

\subsection{Headline speedup}

\begin{table}[H]
\centering
\caption{Phase 2 speedup (mean and best across the sweep), Phase 2 simulator.}
\label{tab:speedup}
\begin{tabular}{llrrr}
\toprule
benchmark    & policy   & $n$ & avg speedup & best speedup \\
\midrule
libquantum   & off      & 8 & 1.0000 & 1.0000 \\
libquantum   & nextline & 8 & 5.8301 & \textbf{5.8327} \\
libquantum   & adaptive & 8 & 3.8808 & \textbf{5.0245} \\
\midrule
hmmer        & off      & 8 & 1.0000 & 1.0000 \\
hmmer        & nextline & 8 & 1.0058 & 1.0085 \\
hmmer        & adaptive & 8 & 1.0004 & 1.0009 \\
\midrule
dealII       & off      & 8 & 1.0000 & 1.0000 \\
dealII       & nextline & 8 & 1.0158 & 1.0182 \\
dealII       & adaptive & 8 & 1.0107 & 1.0114 \\
\bottomrule
\end{tabular}
\end{table}

Adaptive recovers most of next-line's \texttt{libquantum} win
($5.02\times$ best vs.\ $5.83\times$). On \texttt{hmmer} and \texttt{dealII}
adaptive's speedup is slightly lower than next-line in absolute terms; the
difference is small in cycles and is more than recovered when one accounts
for prefetch accuracy (next subsection).

\subsection{Accuracy and coverage --- the bandwidth lens}
Speedup is not the whole story. A prefetch that is \emph{issued but never
consumed} costs L2 bandwidth, hurts power, and contends with demand
traffic. The accuracy metric (Table~\ref{tab:accuracy}) shows where the
adaptive policy earns its keep:

\begin{table}[H]
\centering
\caption{Adaptive ($D=2$) vs.\ next-line: prefetch accuracy.}
\label{tab:accuracy}
\begin{tabular}{lrr}
\toprule
benchmark   & nextline accuracy & adaptive ($D=2$) accuracy \\
\midrule
libquantum  & $99.4\%$ & $99.4\%$ \\
hmmer       & $93.4\%$ & $93.2\%$ \\
dealII      & $78.4\%$ & $\mathbf{83.3\%}$ \\
\bottomrule
\end{tabular}
\end{table}

At \texttt{max\_prefetch\_depth=2}, adaptive matches or beats next-line on
accuracy across all three benchmarks. The win on \texttt{dealII}
($+5\,\text{pp}$) is the cleanest example of the policy doing what
the paper claims: \emph{not} prefetching when the histogram says streams
in the current epoch are too short to justify it.

\subsection{Adaptive depth-vs-accuracy trade-off}
Table~\ref{tab:depth-tradeoff} breaks the adaptive runs out by the
\texttt{max\_prefetch\_depth} knob. At $D=8$ adaptive becomes more
aggressive, gaining coverage but losing accuracy on irregular workloads.
At $D=2$ accuracy is highest. The user can dial this knob without
re-tuning the policy --- the histogram learning is invariant.

\begin{table}[H]
\centering
\caption{Adaptive policy: per-benchmark accuracy (mean) by max\_prefetch\_depth.}
\label{tab:depth-tradeoff}
\begin{tabular}{lrr}
\toprule
benchmark   & $D=2$ accuracy & $D=8$ accuracy \\
\midrule
libquantum  & $99.4\%$ & $99.4\%$ \\
hmmer       & $93.2\%$ & $69.7\%$ \\
dealII      & $82.3\%$ & $54.2\%$ \\
\bottomrule
\end{tabular}
\end{table}

%─────────────────────────────────────────────────────────────────────────────
\section{Discussion}
%─────────────────────────────────────────────────────────────────────────────

\subsection{Two phases, one story}
Phase~1 and Phase~2 measure different things --- Phase~1 reports raw
L1D miss reduction and L2 bandwidth, Phase~2 reports cycle-modeled
speedup --- but they agree on the qualitative picture:

\begin{enumerate}
  \item \texttt{libquantum} is dominated by long sequential streams;
        prefetching is the right call at any reasonable depth.
        Phase~1 measures this as a $43\times$ miss-count collapse;
        Phase~2 measures it as a $\sim$5$\times$ cycle speedup.
  \item \texttt{hmmer} and \texttt{dealII} have access patterns that
        are mostly orthogonal to consecutive-line prefetching.
        Phase~1 quantifies the failure mode as wasted bandwidth
        ($9\times$ L2 traffic at depth~$8$ on \texttt{hmmer} for
        $2\%$ miss reduction). Phase~2 shows that an adaptive policy
        \emph{at depth $2$} avoids this failure mode while preserving
        the small available speedup.
  \item Aggressive uniform prefetching (Phase~1 at high depth) is
        unambiguously bad on irregular workloads. Adaptive prefetching
        (Phase~2) is dominated by next-line in raw cycles only because
        our cycle model rewards rescued misses without penalizing wasted
        L2 bandwidth. On a hardware platform where wasted prefetches
        cost real DRAM cycles, contention, or power, the accuracy
        advantage in Table~\ref{tab:accuracy} would translate into a
        cycle-time advantage as well.
\end{enumerate}

\subsection{What we learned}
The three bug fixes in Section~\ref{sec:bugs} were instructive: a
plausible-looking adaptive policy can fail silently in two distinct
ways (the histogram-cap cliff, the bootstrap dead-zone), and these
failures present the same way externally --- ``the policy is too
conservative.'' The diagnostic that finally cracked it was inspecting
per-configuration prefetch counts and noticing that adaptive issued
$\sim 4\times$ fewer prefetches than next-line on \texttt{libquantum}
despite an even higher per-prefetch accuracy. The per-config CSV the
sweeper emits (cycles, accuracy, coverage) was indispensable; the
plots alone would have shown the symptom but not the cause.

\subsection{Limitations and future work}
\begin{itemize}
  \item Our simulator is single-level associative on each cache; we did
        not explore higher associativities for L1D/L2. Set-associativity
        primarily affects baseline miss rate and thus is orthogonal to
        the prefetcher comparison, but a richer sweep would quantify
        that orthogonality.
  \item Our latency model treats all L2 misses as $80$ refs, ignoring
        bank conflicts, queueing, and bandwidth contention. The accuracy
        advantage of adaptive (Table~\ref{tab:accuracy}) would translate
        into measurable cycle wins on a more realistic memory model.
  \item We did not model write-back traffic or coherency, both of which
        interact with prefetch effectiveness on multi-core systems.
  \item We omitted strided and indirect prefetchers from the comparison;
        the literature documents large gains for strided detection on
        \texttt{dealII}-class workloads.
\end{itemize}

%─────────────────────────────────────────────────────────────────────────────
\section{Reproducibility}
%─────────────────────────────────────────────────────────────────────────────
The full experimental pipeline is one command:

\begin{lstlisting}[language=bash]
export PIN_ROOT=/path/to/intel-pin
cd ali && bash run.sh
\end{lstlisting}

\noindent
\texttt{run.sh} verifies \texttt{PIN\_ROOT}, builds the Pintool if missing,
sweeps the 72 configurations described in Section~\ref{sec:methodology},
and generates 17 plot files in \texttt{ali/plots/}. The sweep is
resumable: killing it at any time and re-running picks up from the
partial CSV. The Phase~1 sweep is similarly invoked with
\texttt{cd aengus \&\& bash sweep.sh}. Both phases share the
\texttt{benchmarks/} directory and a Python virtual environment at
\texttt{.venv/}.

%─────────────────────────────────────────────────────────────────────────────
\section{Conclusion}
%─────────────────────────────────────────────────────────────────────────────
We have implemented and evaluated two stream-buffer prefetchers on a
shared Pin-based simulator. The fixed-depth Jouppi design (Phase~1)
demonstrates the workload-sensitivity of stream prefetching:
\texttt{libquantum} benefits massively, \texttt{hmmer} and
\texttt{dealII} mostly do not. The adaptive Palacharla--Kessler design
(Phase~2) reproduces that finding under a richer cycle model and adds a
tunable bandwidth-vs-coverage knob: at \texttt{max\_prefetch\_depth=2}
the adaptive policy attains higher prefetch accuracy than next-line on
\texttt{dealII} ($83\%$ vs.\ $78\%$); at maximum depth it captures
$5.0\times$ speedup on \texttt{libquantum} (vs.\ next-line's $5.8\times$).
The three implementation bugs we identified and fixed
(Section~\ref{sec:bugs}) underscore that adaptive prefetcher correctness
is subtle --- two plausible-looking design choices can each silently
disable the policy on its target workload --- and that per-configuration
accuracy/coverage logging is essential to surface those failures.

%─────────────────────────────────────────────────────────────────────────────
\begin{thebibliography}{9}
\bibitem{jouppi1990improving}
N.~P.~Jouppi, ``Improving Direct-Mapped Cache Performance by the Addition
of a Small Fully-Associative Cache and Prefetch Buffers,''
\textit{17th Annual International Symposium on Computer Architecture (ISCA)},
1990.

\bibitem{palacharla1994evaluating}
S.~Palacharla and R.~E.~Kessler, ``Evaluating Stream Buffers as a Secondary
Cache Replacement,''
\textit{21st Annual International Symposium on Computer Architecture (ISCA)},
1994.
\end{thebibliography}

\end{document}
